\section{The Positron As Holonomy Sign}

A sign is not read on an isolated segment. A segment can carry transport, a
local comparison, a change of label, or a bounded residue, but it has not yet
returned to the reference that would make orientation accountable. The grammar
therefore forbids charge from appearing merely because a path has been
walked. Charge appears when the path closes and the carried orientation
returns with a sign the quotient cannot erase.

This is the holonomy reading. Transport a distinction around a comparison.
If the local pieces can all be identified and the return also identifies, the
closed reading is trivial. If the local pieces look admissible but the return
carries a signed residue, the loop has measured orientation. The sign is not
attached to a substance after the fact. It is the record of how a
distinction came home.

Open paths therefore read differently from closed loops. An open path may
carry a phase, a baseline, a local successor count, or a direction of
refinement. It may be essential for later comparison. But until the path is
closed, the grammar has no final return against which to read the sign. The
open-path charge is zero in this precise sense: not because nothing has
happened, but because the orientation debt has not been returned to the
ledger as a closed residue.

The loop can return with a sign. When it does, the sign belongs to the closed
comparison rather than to any one segment. A local reader may find no field
strength along each branch and still read a phase shift when the branches
recombine. A pair of refinement clocks may each move admissibly around a
boundary and still fail to synchronize when carried through the full cycle.
These examples are familiar in later physical registers, but the
grammatical lesson is smaller and stricter: local triviality does not force
global triviality. A loop may close with a holonomy.

The positron is read at exactly this level. The previous charge grammar
already fixed the native flat baseline: over that comparison the negative
sign is the available charged reading, and the positive sign is forbidden.
To read the positive sign, the grammar must select an oriented or tilted
comparison. The same carried activity is then returned through a loop whose
orientation faces the reader oppositely. If the closed comparison comes home
with the positive sign, the grammar reads positron.

This does not import antimatter as a second stock of objects. It reads the
positive charge as the holonomy sign of a closed comparison. The native
baseline did not contain a hidden positive particle waiting to be uncovered.
It contained no admitted positive sign. The sign appears when the grammar
has selected the tilted comparison and closed the orientation. The positron
is therefore a reading of closure, not an added primitive beside the
electron.

The distinction between open and closed is also why annihilation can serve
only as an illustration. A finite coincidence reader may register two
opposed photons inside an energy window and in opposite directions. Such a
record is strong evidence, in its physical setting, that an electron and a
positron have cancelled. But the coincidence reader uses a sign distinction
already made available by the grammar. It does not create the positive sign.
It illustrates how a later finite ledger can read cancellation once closed
orientation has supplied the two opposed charges.

The Aharonov--Bohm illustration is closer to the holonomy form. Two branches
may travel where the local field reading is silent and still recombine with
a phase difference because the connection has carried memory around the
loop. The grammar does not need that experiment to authorize the positron
reading. It uses the shape as a lantern: a closed comparison can preserve a
global sign that no open segment by itself can spend.

The Sagnac illustration gives the clock face of the same obligation. Two
paths traverse a closed boundary in opposite directions and return with
different tallies. The difference is not a defect in local propagation. It is
the holonomy of the bookkeeping rule. Again, the chapter's claim is
grammatical before it is physical: closure can expose a residue that local
transport did not settle.

Chirality supplies a final caution. Left-oriented and right-oriented
refinements may agree in scalar readings and still refuse to substitute for
one another without cost. The difference is oriented admissibility. A mirror
turn is not automatically the same turn. The positron sign depends on this
kind of oriented comparison. If the grammar allowed a handed turn to be
erased without return, it would erase the condition under which the positive
sign can be read.

The section's rule is therefore simple. Open paths can carry comparison but
not charge. Closed loops can read a holonomy sign. The flat native baseline
reads the negative sign and forbids the positive one there. A tilted or
oriented closure may read the positive sign. The positron is that positive
closed-loop reading. The grammar permits it only where the loop closes,
orients, and returns with a sign that cannot be quotiented away.

The next question is a count. If the native baseline admits the negative
reading and forbids the positive one until a tilted closure is selected, what
does the native ledger count before the tilt? The answer is the chapter's
native asymmetry: matter is counted, antimatter is not.
